\documentclass[tikz,border=10pt]{standalone}
\usepackage{ctex}
\usepackage{amsmath}
\usetikzlibrary{calc, arrows.meta}

\begin{document}
% geo-p9-06-1: 3 类空间距离对比图
% （点-平面距离 / 点-直线距离 / 异面直线距离）
\begin{tikzpicture}[
  x={(1cm,0cm)},
  y={(0.40cm,0.32cm)},
  z={(0cm,1cm)},
  thick,
  >=Stealth,
  line cap=round, line join=round,
]

%% ============================================================
%% 栏 1：点到平面距离
%% ============================================================
\begin{scope}[xshift=0cm]
  \node[font=\bfseries\small] at (1.5, -0.6, 3.8) {(1) 点到平面距离};
  \node[font=\small] at (1.5, -0.6, 3.2)
    {$d=\dfrac{|\vec{AP}\cdot\vec{n}|}{|\vec{n}|}$};

  % 平面 π（浅蓝色）
  \fill[blue!10, draw=blue!40, thin]
    (-0.5, -0.5, 0) -- (3.5, -0.5, 0) -- (3.5, 3.0, 0) -- (-0.5, 3.0, 0) -- cycle;
  \node[blue!60, font=\small] at (3.2, 2.7, 0) {平面 $\pi$};

  % 面内一点 A
  \coordinate (A1) at (1.5, 1.5, 0);
  \fill[blue!70] (A1) circle [radius=1.8pt];
  \node[below right, font=\small] at (A1) {$A$（面内点）};

  % 法向量 n（向上）
  \draw[blue!70, dashed, ->] (A1) -- ++(0, 0, 2.0)
    node[above, font=\small, blue!80] {$\vec{n}$};

  % 空间点 P（红，在平面上方）
  \coordinate (P1) at (2.5, 0.5, 2.2);
  \fill[red] (P1) circle [radius=2pt];
  \node[right, font=\small, red] at (P1) {$P$};

  % AP 向量
  \draw[red!80, ->] (A1) -- (P1) node[midway, right=2pt, font=\small] {$\vec{AP}$};

  % 垂足（P 到平面的垂足）
  \coordinate (Foot1) at (2.5, 0.5, 0);
  \fill[orange] (Foot1) circle [radius=1.5pt];
  \node[below, font=\small, orange!80!black] at (Foot1) {垂足 $H$};

  % 距离 d（P 到平面的铅垂线）
  \draw[red, very thick] (P1) -- (Foot1) node[midway, right=2pt, font=\small, red!80!black] {$d$};

  % 直角标记
  \draw[orange!80!black, thin]
    ($(Foot1) + (0.12, 0, 0)$) -- ++(0, 0, 0.12) -- ++(-0.12, 0, 0);

  \node[font=\small, align=center, text width=4cm] at (1.5, -0.6, -1.3)
    {$A$ 为面内任意点\\$\vec{AP}$ 在 $\vec{n}$ 方向的投影长 $= d$};
\end{scope}

%% ============================================================
%% 栏 2：点到直线距离
%% ============================================================
\begin{scope}[xshift=5.8cm]
  \node[font=\bfseries\small] at (1.5, -0.6, 3.8) {(2) 点到直线距离};
  \node[font=\small] at (1.5, -0.6, 3.2)
    {$d=\sqrt{|\vec{AP}|^2 - \left(\dfrac{\vec{AP}\cdot\vec{l}}{|\vec{l}|}\right)^2}$};

  % 直线 l（黑，水平）
  \coordinate (lA2) at (-0.2, 1.5, 1.0);
  \coordinate (lB2) at (3.8, 1.5, 1.0);
  \draw[black, very thick, ->] (lA2) -- (lB2) node[right, font=\small] {$l$};

  % 方向向量标注
  \draw[gray, ->] (1.8, 1.5, 1.0) -- ++(0.7, 0, 0)
    node[above, font=\small] {$\vec{l}$};

  % 直线上一点 A
  \coordinate (A2) at (0.8, 1.5, 1.0);
  \fill[blue!70] (A2) circle [radius=1.8pt];
  \node[below left, font=\small] at (A2) {$A$（线上点）};

  % 空间点 P（红，在直线旁边）
  \coordinate (P2) at (2.5, -0.2, 2.5);
  \fill[red] (P2) circle [radius=2pt];
  \node[below right, font=\small, red] at (P2) {$P$};

  % AP 向量
  \draw[red!80, ->] (A2) -- (P2) node[midway, left=2pt, font=\small] {$\vec{AP}$};

  % P 到直线的垂足 H（AP 在 l 方向投影终点）
  % 投影向量：proj = (AP · l̂) l̂  ，简化手绘位置
  \coordinate (Foot2) at (2.5, 1.5, 1.0);
  \fill[orange] (Foot2) circle [radius=1.5pt];
  \node[above, font=\small, orange!80!black] at (Foot2) {垂足 $H$};

  % A 到 H 的投影（投影分量）
  \draw[green!60!black, dashed, ->] (A2) -- (Foot2)
    node[midway, above=2pt, font=\small, green!70!black] {投影};

  % H 到 P 的距离（垂线，即所求 d）
  \draw[red, very thick] (Foot2) -- (P2)
    node[midway, right=2pt, font=\small, red!80!black] {$d$};

  % 直角标记
  \draw[orange!80!black, thin]
    ($(Foot2) + (0, 0, 0.12)$) -- ++(0.12, 0, 0) -- ++(0, 0, -0.12);

  \node[font=\small, align=center, text width=4.5cm] at (1.5, -0.6, -1.3)
    {勾股：$|\vec{AH}|$ = 投影长\\$d^2 = |\vec{AP}|^2 - |\vec{AH}|^2$};
\end{scope}

%% ============================================================
%% 栏 3：异面直线之间的距离
%% ============================================================
\begin{scope}[xshift=11.8cm]
  \node[font=\bfseries\small] at (1.5, -0.6, 3.8) {(3) 异面直线距离};
  \node[font=\small] at (1.5, -0.6, 3.2)
    {$d=\dfrac{|\vec{AB}\cdot\vec{n}|}{|\vec{n}|}$};

  % 直线 l1（蓝，偏低）
  \coordinate (l1A) at (-0.5, 0.5, 0.5);
  \coordinate (l1B) at (3.5, 0.5, 0.5);
  \draw[blue!70, very thick, ->] (l1A) -- (l1B) node[right, font=\small, blue!80] {$l_1$};

  % 直线 l2（绿，偏高，方向不同——异面）
  \coordinate (l2A) at (0.5, -0.5, 2.5);
  \coordinate (l2B) at (2.5, 3.2, 2.5);
  \draw[green!70!black, very thick, ->] (l2A) -- (l2B) node[right, font=\small, green!70!black] {$l_2$};

  % l1 上点 A，l2 上点 B
  \coordinate (A3) at (0.5, 0.5, 0.5);
  \coordinate (B3) at (1.5, 1.5, 2.5);
  \fill[blue!70] (A3) circle [radius=2pt];
  \fill[green!70!black] (B3) circle [radius=2pt];
  \node[below left, font=\small] at (A3) {$A$};
  \node[right, font=\small] at (B3) {$B$};

  % AB 向量
  \draw[purple, dashed, ->] (A3) -- (B3)
    node[midway, left=2pt, font=\small, purple] {$\vec{AB}$};

  % 公共法向量 n（⊥ l1 且 ⊥ l2）——画在两线中间位置
  \coordinate (Nmid) at (2.0, 0.5, 1.5);
  \draw[red!70!black, very thick, ->] (Nmid) -- ++(0, 0, 1.3)
    node[above, font=\small, red!80!black] {$\vec{n}$（公共法向量）};
  \node[red!60, font=\tiny, right=2pt] at ($(Nmid) + (0,0,0.6)$)
    {$\vec{n}\perp l_1,\, \vec{n}\perp l_2$};

  % 两线间公垂线（最短距离）
  \coordinate (H3A) at (1.5, 0.5, 0.5);   % l1 上的公垂足
  \coordinate (H3B) at (1.5, 0.5, 2.5);   % l2 上的公垂足（简化，竖直）
  \draw[red, very thick, <->] (H3A) -- (H3B)
    node[midway, right=2pt, font=\small, red!80!black] {$d$（最短距离）};

  \node[font=\small, align=center, text width=4.5cm] at (1.5, -0.6, -1.3)
    {$A\in l_1$，$B\in l_2$\\$\vec{n}=\vec{l_1}\times\vec{l_2}$（叉积）或解方程\\$d = |\vec{AB}|$ 在 $\vec{n}$ 方向的投影};
\end{scope}

\end{tikzpicture}
\end{document}
